\section{Discussion}\label{sec:discussion}

Prediction from survey-estimated curves requires naming the target.
Exchangeability of observed curves yields a finite-sample guarantee for a
further survey estimate, and nothing more. Reaching the population
distribution requires either a sampling-error bound, which is shape-free and
wide, or a scale correction, which is efficient and rests on a distributional
assumption that coordinate variances cannot verify.

Where the correction is reachable is governed by
Proposition~\ref{prop:coupling}, and the answer is narrower than a
population-count threshold suggests. The quantity that makes a correction worth
attempting is the quantity that makes its result hard to certify, because the
diagnostic is evaluated on the scale the correction itself shrinks. A
feasibility region drawn as ``enough populations, and enough noise'' is a
diagonal band, and the reported floor of $94$ is its intercept at a design
share of zero, where nothing needs removing.

The survey evidence is consistent throughout. National estimates sit in the
unnecessary regime. Regional estimates within the same survey generate real
sampling noise, and mostly fail the reliability requirement at the very design
shares that make them candidates. Whether the few surviving configurations
activate depends on the variance estimator, which is not a reassuring place for
a feasibility result to rest.

\paragraph{What we do not claim.} No universal width-optimality for any
latent-target band; the non-identification argument motivates assumptions
rather than certifying a radius. No demonstrated improvement over existing
noisy-calibration or small-area methods, since no matched comparison was run.
Activation of a selector is not evidence that~(S) holds: the gates read scale
precision, and a procedure that abstains has not thereby proved its
approximation accurate anywhere else. Coverage measured against survey
estimates does not validate coverage for latent curves, and where the latent
object is unobserved, as in the regional application, that check is unavailable
in principle rather than merely absent.

\paragraph{For practice.} Three levels should be reported separately and never
interchanged: a finite-sample level for the observed target, a latent-target
level with the assumptions and error terms that support it, and an empirical
coverage assessment under a stated simulation design. Alongside them belong the
unit assumed exchangeable, the grid, and the branch a selector returned. Where
a correction is contemplated, \eqref{eq:boundary} gives the population count
required at the design share actually present, which is a more useful planning
quantity than a fixed floor.

\paragraph{Open problems.} Informative latent-target guarantees under weaker
shape restrictions than~(S), since Section~\ref{sec:simulation} shows the cost
of violating it is directional and grows with simultaneity. A version
of~\eqref{eq:boundary} for biased, shrinkage, or structured-covariance variance
estimators, which the present statement excludes and which are exactly the
estimators small-area practice is built on. And design-variance estimation that
does not force a choice between a downward-biased scheme under which a
correction appears reachable and a defensible one under which it does not.
